\chapter{Типы нейронов}
\label{sec:neurontypes}

\section{Нейрон как чёрный ящик}

Допустим, вам дали мешок с восемьюдесятью шестью миллиардами нейронов~\cite{Azevedo2009} и попросили разложить их по кучкам. Как бы вы это сделали?

Инженер, которому дали мешок незнакомых микросхем, не стал бы сортировать их по форме корпуса. Он спросил бы: сколько у детали входов, сколько выходов и что она выдаёт на выход. Нарисуем нейрон так же --- как блок на схеме (рис.~\ref{fig:blackbox}).

\begin{figure}[t]
\centering
\begin{tikzpicture}[>=Stealth,x=1cm,y=1cm]
% the box
\node[draw,very thick,minimum width=2.6cm,minimum height=2.6cm,fill=blue!6] (N) at (0,0) {};
\node at (0,0.55) {\small нейрон};
\node at (0,-0.15) {$\displaystyle u=\sum_i w_i x_i$};
\node at (0,-0.8) {\small $y=\Theta(u-\theta)$};
% inputs
\foreach \k/\y/\lab in {1/1.05/{x_1},2/0.55/{x_2},3/0.05/{x_3}}{
  \draw[->,thick,blue!60!black] (-4.2,\y) -- (-1.3,\y);
  \node[left] at (-4.2,\y) {$\lab$};
  \node[above,blue!60!black] at (-2.1,\y-0.06) {\scriptsize $w_{\k}$};}
\node at (-4.45,-0.4) {$\vdots$};
\node at (-2.1,-0.4) {$\vdots$};
\draw[->,thick,blue!60!black] (-4.2,-1.05) -- (-1.3,-1.05);
\node[left] at (-4.2,-1.05) {$x_n$};
\node[above,blue!60!black] at (-2.1,-1.11) {\scriptsize $w_{n}$};
\draw[decorate,decoration={brace,amplitude=5pt},gray!60!black] (-4.9,-1.2) -- (-4.9,1.2);
\node[left,align=right,gray!40!black] at (-5.1,0) {\small $n\sim10^3$--$10^5$\\[-2pt]\small входов};
% single output wire, then fan-out
\draw[very thick,red!70!black] (1.3,0) -- (3.2,0);
\foreach \x in {1.6,2.2,2.34,2.9}{\draw[thick,red!70!black] (\x,0) -- (\x,0.4);}
\node[above right,font=\tiny\itshape,gray!30!black,inner sep=1pt] at (1.42,0.45) {щёлк\dots\ щёлк-щёлк\dots};
\node[below,red!70!black,align=center] at (2.25,-0.05) {\scriptsize один выход $y$};
\fill[red!70!black] (3.2,0) circle (0.06);
\foreach \y in {1.3,0.65,0,-0.65,-1.3}{
  \draw[->,thick,red!70!black] (3.2,0) -- (5.0,\y);}
\draw[decorate,decoration={brace,amplitude=5pt},gray!60!black] (5.3,1.45) -- (5.3,-1.45);
\node[right,align=left,gray!40!black] at (5.5,0) {\small копии $y$\\[-2pt]\small на $10^3$--$10^4$\\[-2pt]\small других нейронов};
\end{tikzpicture}
\caption{Нейрон глазами инженера. Много входов $x_i$, каждый со своим весом $w_i$; внутри --- взвешенная сумма $u$ и сравнение с порогом $\theta$ ($\Theta$ --- ступенька: $1$, если аргумент положителен, и $0$ иначе); один выход $y$ --- последовательность импульсов, размноженная ветвлением на тысячи получателей.}
\label{fig:blackbox}
\end{figure}

На выходе блока --- не число, а последовательность коротких импульсов напряжения: «щёлк», пауза, «щёлк-щёлк», длинная пауза. Каждый импульс одинаковой высоты, так что вся информация --- в том, \emph{когда} они происходят. Внутри блока, в первом грубом приближении, стоит сумматор и компаратор: входы складываются с весами, и если сумма перевалила за порог, блок выдаёт импульс. В следующей главе мы заменим его моделью, которая честно работает в непрерывном времени.

Выход один, но провод ветвится, и каждая ветка несёт \emph{ту же самую} копию сигнала. В электронике это называется разветвлением по выходу, fan-out. У нейрона коры он порядка нескольких тысяч; логический вентиль в микросхеме обычно нагружен на единицы других вентилей.

\emph{Что} же передаётся по выходному проводу к получателю? Импульс добегает до конца каждой ветки и там превращается в химический сигнал: ветка выбрасывает в узкую щель между клетками щепотку определённого вещества, \emph{нейромедиатора}, которое меняет входной ток получателя. Вот и признак для сортировки: \emph{какой медиатор выбрасывает выход нейрона}.

\section{Один выход --- один тип сигнала}

Сортировка имеет смысл, только если у каждого нейрона тип выхода один. Так ли это? Опыт говорит, что почти так.

\begin{keyblock}
\begin{proposition}[Один нейрон --- один тип выхода]
\label{prop:dale}
У почти каждого нейрона один главный медиатор, и выбрасывает он его во всех ветках своего выхода. Поэтому тип нейрона задаётся его главным медиатором~\cite{Strata1999}.
\end{proposition}
\end{keyblock}

Договоримся обозначать тип нейрона \emph{первыми четырьмя буквами английского названия его главного медиатора}. Так, нейрон, главный медиатор которого глутамат, --- это \nt{GLUT}-нейрон. Все типы собраны в таблице~\ref{tab:types}.

\begin{table}[t]
\centering
\footnotesize
\setlength{\tabcolsep}{3pt}
\renewcommand{\arraystretch}{1.15}
\begin{tabular}{>{\raggedright}p{1.0cm}>{\raggedright}p{2.05cm}>{\raggedright}p{1.75cm}>{\raggedright}p{1.6cm}>{\centering}p{0.7cm}>{\raggedright}p{1.55cm}>{\raggedright}p{1.8cm}>{\raggedright\arraybackslash}p{3.2cm}}
\toprule
Тип & Главный медиатор & Шина & Скорость & Знак & Нейронов в мозге & Выходов у нейрона & Роль в вычислении \\
\midrule
\nt{GLUT} & глутамат & адресная & быстро и медленно & $+$ & $\sim8\cdot10^{10}$ & $\sim10^4$ & главный положительный вес \\
\nt{GABA} & GABA & адресная & быстро и медленно & $-$ & $\sim3\cdot10^{9}$ & $10^3$--$10^4$ & главный отрицательный вес \\
\nt{GLYC} & глицин & адресная & быстро & $-$ & $10^6$--$10^7$ & $10^2$--$10^3$ & отрицательный вес в спинном мозге и стволе; второй ключ для одного из рецепторов глутамата \\
\nt{ACET} & ацетилхолин & обе & быстро и медленно & $\pm$ & $\sim10^6$ & мышца: $10$--$10^3$; мозг: $\sim10^5$ & выход на мышцу; в мозге --- скорость обучения (гипотеза) \\
\nt{DOPA} & дофамин & широковещ. & медленно & $\pm$ & $\sim5\cdot10^{5}$ & $10^5$--$10^6$ & ошибка предсказания награды $\delta$ \\
\nt{SERO} & серотонин & широковещ. & медленно (один тип рецептора быстрый) & $\pm$ & $\sim3\cdot10^{5}$ & $\sim10^5$ & горизонт планирования (гипотеза) \\
\nt{HIST} & гистамин & широковещ. & медленно & $+$ & $\sim6\cdot10^{4}$ & нет оценки & глобальный сигнал «бодрствуй» \\
\nt{NORA} & норадреналин & широковещ. & медленно & $\pm$ & $\sim5\cdot10^{4}$ & $\sim10^5$ & коэффициент усиления (гипотеза) \\
\nt{ADRE} & адреналин & широковещ. & медленно & $\pm$ & $\sim10^4$ & нет оценки & немного нейронов; роль в мозге неясна \\
\nt{ASPA} & аспартат & адресная & быстро & $+$ & нет оценки & нет оценки & слабый положительный вес; роль спорна \\
\bottomrule
\end{tabular}
\caption{Типы нейронов, по убыванию числа нейронов в мозге. \emph{Шина}: адресная --- медиатор выбрасывается в узкую щель к определённому получателю; широковещательная --- от малой группы нейронов-источников на большие области (раздел~\ref{sec:buses}). \emph{Скорость}: быстро --- через ионотропные рецепторы, миллисекунды; медленно --- через метаботропные, сотни миллисекунд и дольше. \emph{Знак} --- обычный; окончательно его задаёт получатель (раздел~\ref{sec:receiver}), и $\pm$ означает, что встречаются оба знака. \emph{Нейронов в мозге} --- сколько нейронов этого типа во всём мозге человека, по порядку величины; всего нейронов около $8.6\cdot10^{10}$~\cite{Azevedo2009}. Почти все \nt{GLUT}-нейроны --- около $7\cdot10^{10}$ --- это мелкие входные нейроны мозжечка; числа \nt{GLUT} и \nt{GABA} получены из этого подсчёта и доли \nt{GABA}-нейронов в коре~\cite{Hendry1987}. Из малочисленных типов прямо подсчитаны только \nt{HIST}~\cite{Airaksinen1991} и главная из двух групп \nt{DOPA}-нейронов~\cite{Pakkenberg1991}, так что для \nt{DOPA} это нижняя граница; для \nt{SERO} --- сумма двух частичных подсчётов~\cite{Baker1990,Baker1991}. Числа для \nt{NORA}, \nt{GLYC}, \nt{ACET} и \nt{ADRE} --- грубые оценки по порядку величины, без прямого подсчёта. \emph{Выходов у нейрона} --- сколько у одного нейрона синаптических окончаний, то есть мест выброса медиатора на ветвях выходного провода (рис.~\ref{fig:blackbox}), по порядку величины. Для широковещательных типов окончания прямо не считали: у \nt{DOPA} число выведено из того, какую долю своей мишени покрывает ветвление одного выходного провода~\cite{Matsuda2009}, у остальных оценки ещё грубее. У \nt{ACET} два случая: нейрон, управляющий мышцей, и широковещательный нейрон в мозге.}
\label{tab:types}
\end{table}

% the pie is a table-type float with a figure caption, so it can never float ahead of Table~\ref{tab:types}
\begin{table}[tb]
\makeatletter\def\@captype{figure}\makeatother
\centering
\definecolor{vizglut}{HTML}{2A78D6}
\definecolor{vizgaba}{HTML}{E34948}
\definecolor{vizrest}{HTML}{8A8986}
\begin{tikzpicture}[>=Stealth]
% pie: GLUT 96.4 %, GABA 3.6 % (13.0 degrees), the rest 0.006 % (a hairline at 0 degrees)
\fill[vizglut] (0,0) -- (13:1.9) arc (13:360:1.9) -- cycle;
\fill[vizgaba] (0,0) -- (0:1.9) arc (0:13:1.9) -- cycle;
\draw[white,line width=1pt] (0,0) -- (13:1.9);
\draw[vizrest,line width=1.2pt] (0,0) -- (0:1.9);
\node[white,align=center,font=\small] at (190:0.95) {\nt{GLUT}\\$96.4\,\%$};
\draw[gray!60!black] (6.5:1.75) -- (2.05,2.1);
\node[anchor=west,font=\small] at (2.05,2.1) {\nt{GABA} $3.6\,\%$};
% zoom box for the remaining types
\draw[densely dashed,vizrest] (1.9,0) -- (3.1,1.65);
\draw[densely dashed,vizrest] (1.9,0) -- (3.1,-2.2);
\draw[vizrest,rounded corners=3pt] (3.1,-2.2) rectangle (10.1,1.65);
\node[anchor=west,font=\scriptsize] at (3.2,1.4) {все остальные вместе: $\approx0.006\,\%$; логарифмическая шкала};
\foreach \code/\lg/\y/\val/\lx in {GLYC/6.477/1.0/{$10^6$--$10^7$}/4.05,ACET/6.0/0.6/{$\sim10^6$}/3.0,DOPA/5.699/0.2/{$\sim5\cdot10^5$}/2.699,SERO/5.477/-0.2/{$\sim3\cdot10^5$}/2.477,HIST/4.778/-0.6/{$\sim6\cdot10^4$}/1.778,NORA/4.699/-1.0/{$\sim5\cdot10^4$}/1.699,ADRE/4.0/-1.4/{$\sim10^4$}/1.0}{
  \node[anchor=east,font=\scriptsize] at (4.35,\y) {\nt{\code}};
  \fill[vizrest] (4.45,\y-0.13) rectangle ({4.45+\lg-3},\y+0.13);
  \node[anchor=west,font=\scriptsize] at ({4.5+\lx},\y) {\val};}
\draw[vizrest,thick] (7.45,1.0) -- (8.45,1.0);
\draw[vizrest,thick] (7.45,0.92) -- (7.45,1.08);
\draw[vizrest,thick] (8.45,0.92) -- (8.45,1.08);
\draw[gray!60!black] (4.45,-1.72) -- (8.45,-1.72);
\foreach \e in {3,4,5,6,7}{
  \draw[gray!60!black] ({4.45+\e-3},-1.72) -- ({4.45+\e-3},-1.66);
  \node[below,font=\scriptsize] at ({4.45+\e-3},-1.72) {$10^{\e}$};}
\end{tikzpicture}
\caption{Доли типов нейронов в мозге по оценкам таблицы~\ref{tab:types}. Круг почти целиком занят \nt{GLUT}; \nt{GABA} --- узкий сектор, а все остальные типы вместе --- волосок на границе. Справа этот волосок в увеличении, по логарифмической шкале числа нейронов; для \nt{GLYC} столбик доходит до середины диапазона $10^6$--$10^7$, отрезок показывает весь диапазон. У \nt{ASPA} оценки нет, и на рисунке его нет.}
\label{fig:typepie}
\end{table}

Насколько неравны эти кучки, видно на рис.~\ref{fig:typepie}: из каждой тысячи нейронов $964$ --- \nt{GLUT}, $36$ --- \nt{GABA}, а на все остальные типы вместе приходится около шести нейронов на сто тысяч.

У GABA название и так состоит из четырёх букв. Вещества, которые почти никогда не бывают главным медиатором, --- нейропептиды, газы, липиды, пурины, --- своего кода не получают; о них --- в приложении~\ref{sec:othermediators}. Слово «почти» в утверждении~\ref{prop:dale} важно. Многие нейроны выбрасывают вместе с главным медиатором вспомогательные вещества~\cite{Yao2023}. Некоторые выбрасывают и второй быстрый медиатор, даже противоположного знака: часть \nt{DOPA}-нейронов выбрасывает ещё и GABA~\cite{Tritsch2012}. А у некоторых нейронов разные медиаторы выходят из разных веток одного и того же выхода~\cite{Zhang2015}. Всё это измерено, так что «один нейрон --- один медиатор» --- правило с исключениями, а не закон. Но как первое деление оно выдерживает проверку: в клеточном атласе мозга мыши, где нейроны разложены по работающим в них генам на пять с лишним тысяч групп, нейроны по-прежнему делятся на крупные классы прежде всего по главному медиатору~\cite{Yao2023}.

У этого правила есть следствие, которое сразу отличает мозг от искусственной нейронной сети. В искусственной сети один и тот же нейрон может иметь положительный вес к одному получателю и отрицательный к другому: веса $w_{ij}$ --- просто числа в матрице. В мозге знак действия в основном определяется медиатором, а медиатор у нейрона один. Значит, если нейрон $j$ возбуждает одного получателя, он возбуждает и всех остальных. Весь \emph{столбец} матрицы весов, выходящий из нейрона $j$, одного знака:
\begin{empheq}[box=\keybox]{equation}
\operatorname{sign} w_{ij} \;=\; s_j \quad\text{для всех получателей } i, \qquad s_j\in\{+1,-1\}.
\label{eq:dalesign}
\end{empheq}
Нейроны делятся на возбуждающие ($s_j=+1$) и тормозящие ($s_j=-1$), и это свойство самого нейрона, а не связи. Если сети нужна отрицательная связь от возбуждающего нейрона, её приходится делать через посредника: возбуждающий нейрон возбуждает тормозящий, а тот тормозит цель: отрицательный вес с задержкой на одно звено.

Медиаторов известно больше сотни, но почти вся работа мозга держится на полудюжине, и они делятся на два совершенно разных класса.

\section{Две шины: адресная и широковещательная}
\label{sec:buses}

В компьютерных сетях есть два способа доставить сообщение. Первый --- \emph{адресный}, unicast: пакет идёт от одного отправителя к определённому получателю, быстро, и больше его никто не видит. Второй --- \emph{широковещательный}, broadcast: отправитель передаёт одно сообщение сразу всем узлам сети. Нейроны пользуются обоими способами, и медиаторы делятся ровно по этому признаку (рис.~\ref{fig:phone_radio}).

\begin{figure}[t]
\centering
\begin{tikzpicture}[>=Stealth,scale=0.95]
% ---- unicast: point-to-point wiring
\node[above] at (2.0,3.3) {\small \textbf{Адресная}: глутамат и GABA};
\foreach \i in {0,...,4}{
  \fill[blue!60!black] (0.2+0.9*\i,2.5) circle (0.1);
  \fill[blue!60!black] (0.2+0.9*\i,0.6) circle (0.1);}
\foreach \a/\b in {0/0,0/1,1/1,1/3,2/2,3/2,3/4,4/4,2/0}{
  \draw[->,blue!60!black] (0.2+0.9*\a,2.38) -- (0.2+0.9*\b,0.72);}
\fill[red!70!black] (1.1,1.55) circle (0.09);
\pic[scale=1.2] at (1.75,1.9) {letter};
\draw[->,red!70!black,thick] (1.1,1.55) -- (0.28,0.7);
\draw[->,red!70!black,thick] (1.1,1.55) -- (1.95,0.72);
\node[align=center,below] at (2.0,0.2) {\small миллиарды нейронов,\\[-2pt]\small у каждого свои адресаты,\\[-2pt]\small время: миллисекунды};
% ---- broadcast: one small source, global targets
\begin{scope}[xshift=7.2cm]
\node[above] at (1.8,3.3) {\small \textbf{Широковещательная}: дофамин, \dots};
\foreach \i in {0,...,8}{\fill[blue!60!black] (-0.2+0.5*\i,2.75) circle (0.08);}
\fill[orange!85!black] (1.8,0.35) ellipse (0.35 and 0.18);
\pic[scale=1.5] at (2.7,0.75) {mast};
\foreach \x in {-0.2,0.3,0.8,1.3,1.8,2.3,2.8,3.3,3.8}{
  \draw[->,orange!85!black] (1.8,0.5) .. controls (1.8,1.3) and (\x,1.6) .. (\x,2.62);}
\node[align=center,below] at (1.8,0.1) {\small тысячи --- сотни тысяч нейронов,\\[-2pt]\small одно сообщение всем,\\[-2pt]\small время: сотни миллисекунд и дольше};
\end{scope}
\end{tikzpicture}
\caption{Два способа передачи. Слева --- адресная шина, похожая на почту с адресом на конверте; красным выделен один нейрон и два его получателя. Справа --- широковещательная, как радиостанция: небольшая группа нейронов-источников, и одно и то же сообщение получают все.}
\label{fig:phone_radio}
\end{figure}

\emph{Адресная шина} --- это глутамат и GABA. Их выбрасывает подавляющее большинство нейронов. Каждый выход ведёт к определённым клеткам, медиатор действует только в узкой щели контакта, и действует быстро: за миллисекунду открываются ионные каналы получателя, за несколько миллисекунд всё кончается. Это шина \emph{данных}: по ней передаётся то, что мозг вычисляет.

\emph{Широковещательная шина} --- дофамин, серотонин, норадреналин и отчасти ацетилхолин. Их выбрасывают крошечные группы нейронов, но выход каждого из них ветвится невероятно широко. Медиатор часто выбрасывается не в узкую щель, а в межклеточное пространство, где он расплывается и действует на всех, кто рядом. Действует медленно: не открывает каналы напрямую, а запускает внутри получателя цепочку химических реакций. По ней идут не данные, а \emph{управляющие параметры}, общие для больших участков сети, которые меняют режим её работы: коэффициент усиления, скорость обучения, сигнал «это было хорошо». Поэтому такие медиаторы называют \emph{модуляторами}. Долго считалось, что каждый такой канал --- одно число на весь мозг. Современные измерения уточнили картину: канал --- не один провод, а небольшой жгут параллельных линий. Нейроны-источники одного медиатора делятся на подсистемы, у которых свои получатели и своё сообщение, а сколько медиатора выбросить на месте, подстраивают ещё и местные сигналы~\cite{Ren2018,Poe2020}. Число таких линий --- единицы или десятки, а не миллиарды, так что канал остаётся широковещательным.

Если вам нужна одна картинка из этой главы --- вот она. Мозг --- огромная сеть с адресной передачей данных, над которой висит несколько широковещательных каналов с управляющими сигналами. Программист узнает знакомое устройство: вычисление плюс небольшой набор управляющих переменных, каждая из которых общая для большого участка сети.

\section{\nt{GLUT}: положительный вес}

Глутамат --- главный возбуждающий медиатор мозга. Когда выход нейрона выбрасывает глутамат, у получателя открываются каналы, через которые внутрь входит натрий, напряжение на мембране получателя растёт, и его шанс выдать собственный импульс увеличивается. На языке рис.~\ref{fig:blackbox} это вход с \emph{положительным} весом, $w_i>0$.

Кто выдаёт глутамат? Почти все, кто передаёт данные на большие расстояния: от трёх четвертей до четырёх пятых нейронов коры и большинство нейронов, связывающих разные области мозга. Сеть мозга --- это в первую очередь сеть глутаматных связей.

Одна инженерная деталь. После выброса концентрация глутамата в щели взлетает в тысячи раз, примерно до миллимоля, и спадает с постоянной времени около миллисекунды~\cite{Clements1992}: глутамат расплывается из щели, а специальные белки-насосы откачивают его обратно в клетки. Зачем? Затем же, зачем в линии связи после каждого символа сигнал возвращают к нулю. Если медиатор не убрать, следующий импульс придёт на «занятую» линию, и импульсы с интервалом в несколько миллисекунд сольются.

\section{\nt{GABA}: отрицательный вес}

Представьте сеть, в которой все веса положительны. Если в среднем каждый импульс порождает больше одного следующего импульса, активность растёт экспоненциально и через несколько шагов охватывает всю сеть. Электронщик узнает петлю положительной обратной связи с усилением больше единицы: схема уходит в насыщение. Чтобы этого не случилось, нужны отрицательные веса. Главный их источник --- гамма-аминомасляная кислота, сокращённо GABA.

GABA открывает у получателя каналы, пропускающие ионы хлора. Равновесный потенциал хлора лежит около потенциала покоя или немного ниже, поэтому открытые хлорные каналы удерживают мембрану около покоя и не дают поднять напряжение до порога. Это вход с \emph{отрицательным} весом, $w_i<0$. GABA-нейроны --- от пятой части до четверти нейронов коры~\cite{Hendry1987}. Их выходы короткие, и тормозят они своих ближайших соседей.

Торможение в мозге делает гораздо больше, чем просто вычитание. Оно работает как отрицательная обратная связь, которая держит всю сеть в рабочей точке. Считается, что в нормально работающей коре возбуждающий и тормозящий токи, приходящие на нейрон, почти точно уравновешивают друг друга: такой режим предсказан теорией~\cite{vanVreeswijk1996} и виден в измерениях токов в живой коре~\cite{Haider2006}, и нейрон всё время сидит около порога. Схема, стабилизированная сильной отрицательной обратной связью около неустойчивой точки, быстро и чутко реагирует на малые сигналы --- старый инженерный приём.

\section{\nt{DOPA}: сигнал ошибки}

Первый широковещательный канал --- дофамин. Дофаминовых нейронов у человека порядка полумиллиона, примерно один на сто семьдесят тысяч; столько подсчитано в главной из двух их групп~\cite{Pakkenberg1991}, так что это нижняя граница. Их выходы расходятся к тем областям, которые выбирают действия.

Что передаёт этот канал? Ответ дают опыты, в которых регистрируют импульсы дофаминовых нейронов у животного, получающего награду~\cite{Schultz1997}. Животному время от времени неожиданно дают каплю сока, и дофаминовые нейроны отвечают на неё короткой пачкой импульсов. Потом перед соком начинают зажигать лампочку, всегда за секунду до сока. Когда животное выучивает эту связь, нейроны начинают отвечать пачкой на \emph{лампочку}, а на сам сок, пришедший точно в срок, не отвечают совсем. А если после лампочки сок не дают, то в момент, когда он должен был прийти, нейроны \emph{замолкают}, и их частота падает ниже обычной.

Правило, которое объясняет все три случая (рис.~\ref{fig:rpe}): нейроны сообщают не о том, насколько хорошо, а о том, насколько \emph{лучше или хуже, чем ожидалось}.

\begin{figure}[t]
\centering
\begin{tikzpicture}[>=Stealth,x=3.4cm,y=1cm]
\foreach \y/\lab in {2.8/{неожиданный сок},1.4/{сок после лампочки},0/{лампочка, но сока нет}}{
  \draw[gray!70!black] (0,\y) -- (3,\y);
  \draw[densely dotted,gray!60!black] (0,\y+0.3) -- (3,\y+0.3);
  \node[left,align=right,font=\small] at (-0.03,\y+0.35) {\lab};}
% row 1: juice without a cue -> burst at the juice
\draw[very thick,orange!85!black] plot[domain=0:3,samples=120] (\x,{2.8+0.3+0.75*exp(-((\x-2.08)/0.07)^2)});
\pic[scale=0.8] at (2,2.58) {juice};
% row 2: cue then juice -> burst at the cue, nothing at the juice
\draw[very thick,orange!85!black] plot[domain=0:3,samples=120] (\x,{1.4+0.3+0.75*exp(-((\x-1.08)/0.07)^2)});
\pic[scale=0.85] at (1,1.15) {lamp}; \pic[scale=0.85] at (2,1.15) {juice};
% row 3: cue, juice omitted -> burst at the cue, dip when the juice was due
\draw[very thick,orange!85!black] plot[domain=0:3,samples=120] (\x,{0.3+0.75*exp(-((\x-1.08)/0.07)^2)-0.28*exp(-((\x-2.1)/0.12)^2)});
\pic[scale=0.85] at (1,-0.25) {lamp}; \pic[scale=0.85] at (2,-0.25) {nojuice};
\node[right,font=\scriptsize] at (2.36,0.1) {провал};
\node[right,font=\scriptsize] at (2.13,3.75) {сюрприз!};
\node[left,font=\scriptsize] at (0.97,2.25) {всплеск на лампочку};
\node[above,font=\scriptsize] at (2,1.75) {ничего};
\draw[->,gray!70!black] (0,-0.75) -- (3.05,-0.75) node[right,black,font=\small] {$t$};
\draw[gray!60!black] (1,-0.8) -- (1,-0.7) node[below=2pt,font=\scriptsize] {лампочка, $1\,\text{с}$};
\draw[gray!60!black] (2,-0.8) -- (2,-0.7) node[below=2pt,font=\scriptsize] {сок, $2\,\text{с}$};
\end{tikzpicture}
\caption{Частота импульсов \nt{DOPA}-нейронов в трёх опытах (схематично, по~\cite{Schultz1997}). Пунктир --- ровная фоновая частота. Лампочка --- сигнал, капля --- сок, перечёркнутая капля --- сок, который обещали, но не дали. Отклик --- ошибка предсказания~\eqref{eq:rpe}.}
\label{fig:rpe}
\end{figure}

Программист, знакомый с обучением с подкреплением~\cite{SuttonBarto2018}, узнает эту величину сразу. Агент, который учится выбирать действия, хранит оценку $V(s)$ --- сколько награды он ожидает, находясь в состоянии $s$. После каждого шага он сравнивает ожидание с тем, что получил:
\begin{empheq}[box=\keybox]{equation}
\delta_t \;=\; r_t \;+\; \gamma\,V(s_{t+1}) \;-\; V(s_t) ,
\label{eq:rpe}
\end{empheq}
и поправляет оценку: $V(s_t)\leftarrow V(s_t)+\alpha\,\delta_t$. Здесь $r_t$ --- полученная награда, $\gamma<1$ --- коэффициент дисконтирования (насколько будущая награда дешевле немедленной), $\alpha$ --- скорость обучения. Величина $\delta_t$ называется \emph{ошибкой временных разностей}.

Она ведёт себя в точности как дофаминовые нейроны. Неожиданный сок: $V$ ещё ноль, $r>0$, $\delta>0$. Выученный сок: награду предсказала лампочка, $V(s_t)$ перед соком уже равна $r$, и $\delta=0$. Сок не пришёл: $V(s_t)=r$, а $r_t=0$, $\delta<0$. А всплеск переезжает на лампочку, потому что именно в момент лампочки ожидание скачком выросло: $\gamma V(s_{t+1})-V(s_t)>0$. Шаги $t,t+1,\dots$ здесь --- условность алгоритма: в организме часов нет, и ту же величину в непрерывном времени мы получим в главе~\ref{sec:dofa}.

Значит, дофамин --- это широковещательная рассылка скалярного сигнала ошибки $\delta$ всем, кто должен по нему учиться. В первом приближении --- один провод на весь мозг и одно число в каждый момент времени; насколько этот провод на самом деле жгут, разбирается в главе~\ref{sec:dofaalt}.

\section{Остальные широковещательные каналы}

Для остальных модуляторов есть лишь рабочие гипотезы, которые переводят их роли на тот же язык глобальных параметров~\cite{Doya2002}. Держите их именно как гипотезы.

\emph{\nt{NORA}, норадреналин: коэффициент усиления.} Норадреналиновых нейронов ещё меньше, чем дофаминовых, по грубой оценке несколько десятков тысяч, а их выходы расходятся практически по всему мозгу. Они резко активируются, когда происходит что-то неожиданное или важное. Норадреналин меняет крутизну передаточной характеристики нейронов-получателей: слабые входы становятся слабее, сильные --- сильнее. Измеряется это так: фоновые импульсы получателя подавляются сильнее, чем ответ на вход, и отношение сигнала к шуму растёт~\cite{Foote1975NE}. В простой модели это описывают одним числом: если нейрон отвечает частотой $f=F\bigl(g\cdot u\bigr)$ на входной ток $u$, то норадреналин увеличивает коэффициент $g$~\cite{AstonJones2005} (подробнее --- в главе~\ref{sec:nora}). Сеть с большим $g$ работает «контрастнее» и быстрее принимает решения; с маленьким --- «размыто» и больше перебирает варианты, как агент обучения с подкреплением в режиме поиска.

\emph{\nt{ACET}, ацетилхолин: скорость обучения.} Ацетилхолин управляет тем, насколько сеть слушает текущий вход, а насколько --- собственные сохранённые данные. При высоком уровне ацетилхолина входы с органов чувств усиливаются, а внутренние, «вспоминающие» связи ослабляются, и одновременно облегчается изменение весов~\cite{Hasselmo2006}. Грубо говоря, это переключатель «запись/чтение» и вместе с ним --- скорость обучения $\alpha$.

\emph{\nt{SERO}, серотонин: горизонт планирования.} Что передаёт серотонин, понятно хуже всего. Одна из гипотез связывает его с коэффициентом дисконтирования $\gamma$ из~\eqref{eq:rpe}: высокий уровень серотонина соответствует готовности подождать большую награду, а не хватать маленькую сейчас. Серотониновые нейроны работают ровно и медленно, несколько импульсов в секунду в бодрствовании, и почти замолкают в одной из фаз сна~\cite{JacobsAzmitia1992}.

Все шесть медиаторов собраны в таблице~\ref{tab:transmitters}.

\begin{table}[t]
\centering
\small
\begin{tabular}{>{\raggedright}p{2.4cm}>{\raggedright}p{3.0cm}>{\raggedright}p{2.0cm}>{\raggedright\arraybackslash}p{5.2cm}}
\toprule
Тип нейрона & Доля нейронов & Шина & Роль в вычислении \\
\midrule
\nt{GLUT} (глутамат) & большинство & адресная & положительный вес, $w>0$ \\
\nt{GABA} (GABA) & 20--25\% в коре & адресная & отрицательный вес, $w<0$; стабилизация \\
\nt{DOPA} (дофамин) & $\sim 6\cdot10^{-6}$ & широковещ. & ошибка предсказания $\delta$ \\
\nt{NORA} (норадреналин) & $\sim 6\cdot10^{-7}$ & широковещ. & коэффициент усиления $g$ (гипотеза) \\
\nt{ACET} (ацетилхолин) & малая & обе & скорость обучения $\alpha$; «запись/чтение» (гипотеза) \\
\nt{SERO} (серотонин) & $\sim 3\cdot10^{-6}$ & широковещ. & дисконтирование $\gamma$ (гипотеза) \\
\bottomrule
\end{tabular}
\caption{Главные медиаторы мозга на языке вычислений. Правый столбец --- сильное упрощение: надёжен для глутамата, GABA и дофамина, для остальных это гипотезы.}
\label{tab:transmitters}
\end{table}

\section{Знак задаёт получатель}
\label{sec:receiver}

Я немного вас обманул, говоря: глутамат --- положительный вес, GABA --- отрицательный. Ни одна молекула не несёт в себе знака. Молекула --- это просто ключ. Что произойдёт, когда её поймают, решает \emph{замок} --- рецептор на клетке-получателе.

Лучший пример --- дофамин. У него есть рецепторы двух семейств~\cite{Beaulieu2011}. На одних он облегчает возбуждение клетки, на других --- затрудняет. В области, которая выбирает действия, одни нейроны несут рецепторы первого типа, другие --- второго, и один и тот же дофаминовый сигнал одновременно усиливает одну группу и ослабляет другую. Это как один широковещательный пакет, который разные узлы сети интерпретируют по-разному.

\begin{keyblock}
\begin{proposition}[Смысл сообщения задаёт получатель]
\label{prop:receptor}
Знак и скорость действия медиатора определяются не самим медиатором, а рецептором, с которым он связывается. Рецепторы бывают двух родов. \emph{Ионотропные} сами являются ионными каналами и открываются за доли миллисекунды: это прямое управление током, как затвор транзистора. \emph{Метаботропные} запускают внутри клетки цепочку химических реакций и действуют за сотни миллисекунд и дольше: это скорее запись в регистр настроек, которая меняет последующую работу клетки. Адресная шина говорит в основном через первые, широковещательная --- в основном через вторые.
\end{proposition}
\end{keyblock}

Почему тогда вообще можно сказать «глутамат возбуждает»? Потому что почти все рецепторы глутамата, которые встречаются на практике, возбуждающие, а почти все рецепторы GABA --- тормозные. Это не закон природы, а очень надёжное соглашение, и правило~\eqref{eq:dalesign} на нём и держится.

\section{Что унести с собой}

Подавляющее большинство нейронов --- адресная шина данных с весами одного знака на нейрон; ничтожное меньшинство --- широковещательные каналы, передающие на большие участки мозга несколько общих управляющих сигналов. Около двух нейронов из каждых ста тысяч --- и от них зависит, как учится и в каком режиме работает вся остальная сеть. Для инженера это не удивительно: в любом компьютере регистров управления несравнимо меньше, чем ячеек памяти, и всё же без них машина не работает.

В следующей главе мы проверим, что может вычислить сеть из одних \nt{GLUT}-нейронов, самого многочисленного типа, без всяких часов.

\section{Задачи}

\begin{exercise}[\lvlA]
\label{ex:01:1}
\emph{Кучки и провода.} По таблице~\ref{tab:types} (середина диапазона по логарифмической шкале; \nt{ACET} --- $10^5$ выходов): (а)~если нейрон --- человек, сколько населений Земли составят \nt{GLUT}-нейроны и город какого размера --- все широковещательные? (б)~Во сколько раз доля окончаний \nt{DOPA}, \nt{SERO}, \nt{NORA}, \nt{ACET} больше доли их нейронов?
\end{exercise}

\begin{exercise}[\lvlA]
\label{ex:01:2}
\emph{Возврат к нулю.} Глутамат в щели спадает как $c_0e^{-t/\tau_c}$, $c_0=1\mM$, $\tau_c=1\ms$; получатель замечает концентрацию выше $10$~мкМ. (а)~Сколько времени линия «занята» после одного выброса? (б)~Насосы частично заблокированы, $\tau_c=10\ms$. До какой частоты выбросов линия ещё успевает освободиться?
\end{exercise}

\begin{exercise}[\lvlB]
\label{ex:01:3}
\emph{Куда переезжает всплеск.} В опыте рис.~\ref{fig:rpe} после лампочки идут состояния $s_0\to\dots\to s_{K-1}$ с шагом $\Delta$, $T=K\Delta$; награда $r$ приходит на последнем переходе, момент лампочки непредсказуем. Найдите выученные $V(s_k)$ и ответ $\delta$ на лампочку. Положив $\gamma=e^{-\Delta/\tau_\gamma}$, найдите $\tau_\gamma$ при $T=1$~с, $\Delta=0.1$~с, $\gamma=0.95$.
\end{exercise}

\begin{exercise}[\lvlB]
\label{ex:01:4}
\emph{Моделирование: обучение по ошибке.} Промоделируйте цепочку задачи~\ref{ex:01:3} при $K=10$, $\gamma=0.95$, $r=1$ с правилом $V(s_t)\leftarrow V(s_t)+\alpha\delta_t$. В каком испытании $n^*$ ответ на лампочку впервые обгоняет ответ на сок при $\alpha=0.05$, $0.1$, $0.2$, $0.5$? Как $n^*$ зависит от $\alpha$?
\end{exercise}

\begin{exercise}[\lvlB]
\label{ex:01:5}
\emph{Проектирование: рассылка одного числа.} Число $\delta$ надо доставить всем $10^{10}$ нейронам; каждый получатель, включая ретрансляторы, должен получить $10$ окончаний от предыдущего слоя, звено добавляет $1\ms$. Сколько нейронов и звеньев в самом экономном дереве ретрансляторов при $10^4$ выходах на нейрон (\nt{GLUT}) и при $10^{5.5}$ (\nt{DOPA})?
\end{exercise}

\begin{exercise}[\lvlB]
\label{ex:01:6}
\emph{Почему равновесие чуткое.} В сети $80\,\%$ \nt{GLUT} и $20\,\%$ \nt{GABA}, каждый связан с каждым весами $J_E/N$ и $-J_I/N$; $\tau\dot r=-r+Wr+h$. При $J_E=10$ единственное ненулевое собственное число $W$ равно $0.5$. Найдите отношение тормозящего тока к возбуждающему и на сколько процентов достаточно ослабить торможение, чтобы сеть ушла в лавину.
\end{exercise}

\begin{exercise}[\lvlC]
\label{ex:01:7}
\emph{Исследование: \nt{SERO} --- это $\gamma$?} По задаче~\ref{ex:01:3} ответ \nt{DOPA}-нейронов на лампочку равен $re^{-T/\tau_\gamma}$. Задержки $T=1,\dots,5$~с по $n$ предъявлений, $\ln\delta$ измеряется с разбросом $0.5$. Сколько нужно $n$, чтобы отличить $\tau_\gamma=2$~с от $2.4$~с (уровень $0.05$, мощность $0.8$)? Какой механизм, кроме $\gamma$, дал бы тот же спад с $T$?
\end{exercise}
